\documentclass[12pt,a4paper,sans]{moderncv}

\moderncvcolor{grey}
\moderncvstyle{classic}

\usepackage[utf8]{inputenc}
\usepackage{anyfontsize}
\usepackage[
    scale=0.75,
    noheadfoot,
    nomarginpar,
    top=1.5cm,
    bottom=1.5cm,
    left=1.75cm,
    right=1.75cm,
    % showframe, % uncomment to show how the type block is set on the page
]{geometry}
\recomputelengths

\setlength{\footskip}{37pt}

\firstname{Ozren}
\familyname{Dabić}
\title{Software Engineer}
\email{dabic.ozren@gmail.com}
\homepage{dabico.github.io}
\social[linkedin]{dabico}
\social[github]{dabico}

\nopagenumbers{}

\newcommand{\vspacesmall}{\vspace{0.2cm}}
\newcommand{\cvskillitem}[2]{\cvitem{\textbf{#1}:}{#2}}

\newcommand{\cvlanguageitem}{\cvskillitem}
\newcommand{\cvlanguagedoubleitem}[4]{
    \cvtripleitem{\textbf{#1:}}{\texttt{#2}}{\textbf{#3:}}{\texttt{#4}}{~}{~}
}

\newcommand{\cvreference}[4]{
    \cventry{}{#1}
    {\newline{}#2}
    {\newline{}\texttt{\httplink{#3}}}
    {\newline{}\emaillink{#4}}
    {}
}

\renewcommand*{\bibliographyitemlabel}{[\arabic{enumiv}]}

\begin{document}

\makecvtitle

\section{About}

I specialize in crafting robust, adaptable solutions with meticulous attention
to detail. Although comfortable working as a \textbf{full-stack developer}, my
focus is primarily on \textbf{back-end development} and \textbf{infrastructure maintenance}.
With a dedication to both code quality and maintainability, I strive to create
efficient software that withstands the test of time, offering high configurability
and applicability across diverse contexts. Fueled by curiosity, I continuously
seek to improve my skills and tackle complex challenges head-on. While I am
dedicated and methodical, as demonstrated by my strict adherence to the best
practices of \textbf{software development}, I do not shy away from creative
thinking and exploring unconventional solutions.

%%%%%%%%%%%%%%%%%%%%%%%%%%%%%%%%%%%%%%%%%%%%%%%%%%%%%%%%%%%%%%%%%%%%%%%%%%%%%%%

\section{Experience}

\cventry{Oct 2024 -- Present~~}
{Software Developer}
{\httpslink[JetBrains]{jetbrains.com}}
{Belgrade, Serbia}{}
{
    Working with the \httpslink[Machine Learning Methods in Software Engineering Lab]
    {lp.jetbrains.com/research/ml_methods} (ML4SE) on developing tools and
    services that leverage machine learning to improve the productivity of
    developers.
}

\cventry{Oct 2022 -- Oct 2024~~}
{Research Assistant \& Software Engineer}
{\httpslink[Software Institute]{si.usi.ch}}
{Lugano, Switzerland}{}
{
    Aiding the \httpslink[Software Analytics Research Team]{seart.si.usi.ch}
    (SEART) in conducting research on the usage of machine learning for automation
    of software development processes. This includes creating tools for, and
    conducting empirical studies, while expanding, improving and maintaining
    various web platforms, applications and libraries created while working
    with the group.
}
\cvlistitem{
    Worked on \httpslink[\textbf{GitHub Search}]{seart-ghs.si.usi.ch}, a web
    platform that continiously mines and analyzes open-source projects on GitHub
    while providing users with an interface for dynamically generating a dataset
    by sampling our database based on their specific requirements.
}
\cvlistitem{
    Orchestrated development of \httpslink[\textbf{Data Hub}]{seart-dh.si.usi.ch},
    a web platform that facilitates large-scale mining of source code from open
    source projects hosted on GitHub. Researchers can utilize its intuitive
    interface to define criteria for selecting code and then download the
    generated dataset in order to train deep learning models that work with code.
}
\cvlistitem{
    Created \httpslink[\textbf{labeler}]{github.com/seart-group/labeler}, a web
    application with which researchers can quickly set up and perform collaborative
    manual analysis of text data. It is intended to streamline the process of
    reviewing instances, resolving conflicts and producing a study report.
}
\cvlistitem{
    Implemented \textbf{Java bindings for \httpslink[\texttt{tree-sitter}]
    {github.com/seart-group/java-tree-sitter} and \httpslink[\texttt{CLOC}]
    {github.com/seart-group/jcloc}}. These libraries are used in other projects
    for performing static analysis on mined source code.
}
\cvlistitem{
    Established crucial automated infrastructure, including database backups,
    project dependency updates, code quality checks, and Docker image builds
    and deployments.
}

\cventry{Sept 2019 -- Dec 2019~~}
{Back-end Developer}
{\httpslink[Ex Machina]{exmachina.ch}}
{Lugano, Switzerland}{}
{
    Was part of a small team of developers tasked with extending and maintaining
    a client-management web platform developed for, and used internally by the
    Zürich-based venture capital firm \httpslink[Planven Entrepreneur Ventures]
    {planvenev.com}.
}
\cvlistitem{
    Elevated platform search capabilities by emphasizing dynamically constructed
    SQL queries. This enhancement significantly boosted the flexibility of the
    platform without compromising performance or security.
}
\cvlistitem{
    Implemented a system that facilitates the tracking of industries, allowing
    for dynamic addition, removal, and assignment of companies to specific sectors.
}

%%%%%%%%%%%%%%%%%%%%%%%%%%%%%%%%%%%%%%%%%%%%%%%%%%%%%%%%%%%%%%%%%%%%%%%%%%%%%%%

\section{Publications}

\begin{enumerate}
    \item \textbf{Ozren Dabić}, Emad Aghajani, and Gabriele Bavota.
    Sampling Projects in GitHub for MSR Studies. In 18th IEEE/ACM International
    Conference on Mining Software Repositories, MSR 2021, pages 560-564. IEEE, 2021.
    \item Rosalia Tufano, \textbf{Ozren Dabić}, Antonio Mastropaolo, Matteo Ciniselli,
    and Gabriele Bavota. Code Review Automation: Strengths and Weaknesses of the State
    of the Art. IEEE Transactions on Software Engineering, 50(2):338-353, 2024.
    \item Rosalia Tufano, Antonio Mastropaolo, Federica Pepe, \textbf{Ozren Dabić},
    Massimiliano Di Penta, and Gabriele Bavota. Unveiling ChatGPT's Usage in Open
    Source Projects: A Mining-based Study. In 21st IEEE/ACM International Conference
    on Mining Software Repositories, MSR 2024. IEEE, 2024.
    \item \textbf{Ozren Dabić}, Rosalia Tufano, and Gabriele Bavota.
    SEART Data Hub: Streamlining Large-Scale Source Code Mining and Pre-Processing.
    In 2024 IEEE International Conference on Software Maintenance and Evolution.
    IEEE, 2024.
\end{enumerate}

%%%%%%%%%%%%%%%%%%%%%%%%%%%%%%%%%%%%%%%%%%%%%%%%%%%%%%%%%%%%%%%%%%%%%%%%%%%%%%%

\section{Education}

\cventry{Sep 2020 -- Oct 2022~~}{Master of Science in Software and Data Engineering}
{\newline{}USI - Università della Svizzera italiana}
{Lugano, Switzerland}
{\textit{GPA: 9.66/10}}
{}

\cventry{Sep 2017 -- Jun 2020~~}{Bachelor of Science in Computer Science}
{\newline{}USI - Università della Svizzera italiana}
{Lugano, Switzerland}
{\textit{GPA: 8.25/10}}
{}

%%%%%%%%%%%%%%%%%%%%%%%%%%%%%%%%%%%%%%%%%%%%%%%%%%%%%%%%%%%%%%%%%%%%%%%%%%%%%%%

\section{Academic Supervising}

\cventry{Spring 2024}
{A Platform to Mine Large-Scale and High-Quality Datasets for Code Review Automation}
{Software and Data Engineering Master Thesis}
{Christelle Rossier}{}{}

\cventry{Spring 2023}
{Sampling Open Source Repositories by Mining Code-related Metrics and Repository Topics}
{Computer Science Bachelor Thesis}
{Albert Cerfeda}{}{}

%%%%%%%%%%%%%%%%%%%%%%%%%%%%%%%%%%%%%%%%%%%%%%%%%%%%%%%%%%%%%%%%%%%%%%%%%%%%%%%

\section{Skills}
\cvskillitem{Prog. Lang}{Java, JavaScript, Python, C, Bash, SASS}
\cvskillitem{Framewroks}{Spring Boot, Express.js, Vue, Bootstrap}
\cvskillitem{Databases}{MySQL, PostgreSQL, MongoDB}
\cvskillitem{Automation}{Maven, Flyway, Liquibase, GitHub Actions}
\cvskillitem{Analysis}{Jupyter Notebook, Pandas, Bokeh}
\cvskillitem{Tools}{Git, Docker, Docker Compose, NGINX}

%%%%%%%%%%%%%%%%%%%%%%%%%%%%%%%%%%%%%%%%%%%%%%%%%%%%%%%%%%%%%%%%%%%%%%%%%%%%%%%

\section{Spoken Languages}
\cvlanguagedoubleitem{English}{C2}{Serbian}{C2}
\cvlanguagedoubleitem{German}{B2}{Italian}{A1}

%%%%%%%%%%%%%%%%%%%%%%%%%%%%%%%%%%%%%%%%%%%%%%%%%%%%%%%%%%%%%%%%%%%%%%%%%%%%%%%

\section{References}
\cvreference{Prof. Dr. Gabriele Bavota}
{Head of the Software Analytics Research Team}
{https://www.inf.usi.ch/faculty/bavota}
{gabriele.bavota@usi.ch}

\end{document}
